\chapter{Manual de Usuario}
\label{ch:manual}

Este anexo proporciona las instrucciones necesarias para instalar, configurar y ejecutar el sistema de predicción de activos financieros desarrollado en este trabajo.

\section{Requisitos del Sistema}

\subsection{Software Requerido}
\begin{itemize}
    \item Python 3.11 o superior
    \item pip (gestor de paquetes de Python)
    \item Git (opcional, para clonar el repositorio)
    \item Conexión a Internet (para descargar datos)
    \item No se requieren claves de API (el sentimiento se obtiene de GDELT)
\end{itemize}

\subsection{Hardware Recomendado}
\begin{itemize}
    \item Procesador: Intel Core i5 o equivalente
    \item Memoria RAM: 8 GB mínimo (16 GB recomendado)
    \item Espacio en disco: 1 GB libre
    \item GPU: Opcional, acelera el entrenamiento de modelos DL
\end{itemize}

\section{Instalación}

\subsection{Paso 1: Clonar o Descargar el Proyecto}

Descargue el proyecto o clónelo desde el repositorio:

\begin{verbatim}
git clone <url-del-repositorio>
cd tfg
\end{verbatim}

\subsection{Paso 2: Crear Entorno Virtual (Recomendado)}

\begin{verbatim}
python -m venv .venv
source .venv/bin/activate   # Linux/Mac
.venv\Scripts\activate      # Windows
\end{verbatim}

\subsection{Paso 3: Instalar Dependencias}

\begin{verbatim}
pip install -r requirements.txt
\end{verbatim}

Las principales dependencias incluyen pandas, numpy, yfinance y tensorflow. También se usan scikit-learn, xgboost, statsmodels, matplotlib y ta.

\subsection{Paso 4: Configurar Fuentes de Sentimiento}

El sentimiento específico por activo se obtiene de GDELT y no requiere configurar variables de entorno.

\textbf{Nota:} el descargador usa consultas GDELT abiertas, así que no hay pasos extra de configuración.

\section{Estructura del Proyecto}

\begin{verbatim}
tfg/
+-- data/
|   +-- raw/          # Datos descargados
|   +-- processed/    # Datos preprocesados
+-- src/
|   +-- data/         # Modulos de datos
|   +-- models/       # Implementacion de modelos
|   +-- evaluation/   # Validacion y metricas
|   +-- utils/        # Reproducibilidad y utilidades
+-- scripts/          # Scripts de ejecucion
+-- results/          # Resultados (CSV) y plots
|   +-- plots/        # Figuras generadas
+-- memoria/          # Documento LaTeX
|   +-- Imagenes/     # Figuras incluidas en la memoria
+-- main.py           # Script principal
+-- requirements.txt  # Dependencias
\end{verbatim}

\section{Ejecución del Sistema}

\subsection{Opción 1: Ejecución Completa}

Para ejecutar todo el pipeline (descarga, preprocesado, entrenamiento, evaluación):

\begin{verbatim}
python main.py
\end{verbatim}
Esto genera los CSV en \texttt{results/}, las figuras en \texttt{results/plots/} y las tablas LaTeX en \texttt{memoria/Tablas/}. Además, se copian automáticamente las figuras a \texttt{memoria/Imagenes/} para facilitar la recompilación de la memoria.

\subsection{Opción 2: Ejecución por Etapas}

Si se ejecutan scripts directamente desde la raíz, es necesario que Python pueda encontrar \texttt{src}. En Linux/Mac basta con anteponer \texttt{PYTHONPATH=.} al comando; en Windows puede ejecutarse \texttt{set PYTHONPATH=.} en la consola antes de ejecutar.

\subsubsection{Descargar Datos}
\begin{verbatim}
python -c "from src.data.downloader import DataDownloader; \
           d = DataDownloader(); d.download_all()"
\end{verbatim}
Por defecto descarga el periodo 2020-11-01 a 2025-11-30; si se requiere otro rango, puede ajustarse en \texttt{src/data/downloader.py}.

\subsubsection{Preprocesar Datos}
\begin{verbatim}
PYTHONPATH=. python src/data/preprocessor.py
\end{verbatim}

\subsubsection{Ejecutar Todos los Experimentos}
\begin{verbatim}
PYTHONPATH=. python scripts/run_all_experiments.py
\end{verbatim}
Este script ejecuta los experimentos A-D tal y como se reportan en la memoria.

\subsubsection{Generar Gráficas}
\begin{verbatim}
PYTHONPATH=. python scripts/generate_all_plots.py
PYTHONPATH=. python scripts/generate_summary_plots.py
\end{verbatim}
Las figuras se guardan en \texttt{results/plots/}. Al ejecutar \texttt{main.py} se copian automáticamente a \texttt{memoria/Imagenes/}. Si se ejecutan los scripts por separado, conviene copiar las imágenes a \texttt{memoria/Imagenes/} antes de recompilar el documento.

\subsubsection{Exportar Tablas a LaTeX}

Tras ejecutar los experimentos, las tablas en LaTeX para la memoria pueden generarse con:
\begin{verbatim}
PYTHONPATH=. python scripts/export_tables_to_latex.py
\end{verbatim}
Las tablas se guardan en \texttt{memoria/Tablas/} y se incluyen automáticamente en el documento.

\section{Configuración de Parámetros}

Los principales parámetros configurables se encuentran en los scripts de ejecución:

\begin{itemize}
    \item \texttt{tickers}: Lista de activos a analizar (por defecto: BTC-USD, TSLA, NVDA)
    \item \texttt{horizon}: Horizonte de predicción en días (por defecto: 7)
    \item \texttt{window\_size}: Tamaño de la ventana temporal para DL (por defecto: 30)
    \item \texttt{n\_splits}: Número de iteraciones walk-forward (por defecto: 12)
    \item \texttt{test\_size}: Tamaño del bloque de test por fold (por defecto: 30 días)
    \item \texttt{epochs}: Épocas de entrenamiento para DL (por defecto: 10)
\end{itemize}
Las semillas están fijadas a 42 para facilitar la reproducibilidad y se reinician por fold y modelo en cada entrenamiento. En CPU los resultados son idénticos; en GPU puede haber ligeras variaciones.

\section{Interpretación de Resultados}

Los resultados se guardan en la carpeta \texttt{results/}:

\begin{itemize}
    \item \texttt{experiment\_*\_results.csv}: Métricas de cada experimento
    \item \texttt{directional\_accuracy\_results.csv}: Resultados de precisión direccional
    \item \texttt{diebold\_mariano\_results.csv}: Test de significancia estadística
    \item \texttt{experiment\_c\_ablation\_tests.csv}: Test de Wilcoxon para ablación
    \item \texttt{plots/*.png}: Visualizaciones generadas
\end{itemize}

Las métricas principales son:
\begin{itemize}
    \item \textbf{RMSE}: Error cuadrático medio (en USD)
    \item \textbf{MAE}: Error absoluto medio (en USD)
    \item \textbf{MAPE}: Error porcentual medio (\%) - métrica principal de comparación
    \item \textbf{Directional Accuracy}: Porcentaje de acierto en la dirección del precio (\%)
    \item \textbf{Directional F1}: F1 de la clase ``sube'' (\%)
    \item \textbf{Directional MCC}: Coeficiente de Matthews ([-1, 1])
\end{itemize}

\subsubsection{Calcular Métricas Adicionales}
Para calcular la precisión direccional, F1, MCC y el test de Diebold-Mariano:
\begin{verbatim}
PYTHONPATH=. python scripts/run_additional_metrics.py
\end{verbatim}

\section{Resolución de Problemas}

\subsection{Error de Conexión}
Si falla la descarga de datos, verifique su conexión a Internet y que las APIs estén disponibles.

\subsection{Error de Memoria}
Si el entrenamiento de modelos DL falla por memoria, reduzca el \texttt{batch\_size} o el número de épocas.

\subsection{Advertencias de TensorFlow}
En algunas ejecuciones, TensorFlow puede mostrar avisos internos (por ejemplo, sobre \textit{retracing}). En el código se reduce el nivel de \textit{logging} por defecto para evitar ruido en consola; si aparecen avisos, no afectan a los resultados numéricos. Si se desea ver el log completo para depuración, puede establecerse \texttt{TF\_CPP\_MIN\_LOG\_LEVEL=0} antes de ejecutar.

\section{Contacto y Soporte}

Para dudas o problemas técnicos, se recomienda consultar la documentación del código fuente o contactar con el autor del trabajo.
